% Required packages: \usepackage{booktabs,tabularx}
% Optional for tighter tables: \usepackage{microtype}

\subsection{Token Consumption on MBPP}
\label{sec:mbpp-token-consumption}

We analyze the token consumption of RE-CODE (MetaGPT+EGRCM+PRQE) with
GPT-4o-mini on the 427 completed MBPP problems. In total, the experiment
consumed 3,099,117 tokens, comprising 2,478,098 input tokens (79.96\%) and
621,019 output tokens (20.04\%), or 7,257.89 tokens per completed problem.
At the prices used in our experiment (\$0.15 and \$0.60 per million input
and output tokens, respectively), this corresponds to an estimated cost of
\$0.7443.

\begin{table}[t]
  \centering
  \small
  \caption{Decomposition of token consumption on the 427 completed MBPP
  problems. The reference is the first code-generation call of branch D for
  each problem.}
  \label{tab:mbpp-token-decomposition}
  \begin{tabularx}{\linewidth}{@{}Xrr@{}}
    \toprule
    Component & Tokens & Share (\%) \\
    \midrule
    Reference code generation              &   124,510 &  4.0 \\
    Repeated ensemble generation            & 2,400,020 & 77.4 \\
    Prompt-context inflation                 &    72,115 &  2.3 \\
    EGRCM clarification                      &   202,010 &  6.5 \\
    EGRCM summarization                      &   193,802 &  6.3 \\
    Test-guided repair                       &   106,660 &  3.4 \\
    \midrule
    Total                                    & 3,099,117 & 100.0 \\
    \bottomrule
  \end{tabularx}
\end{table}

\begin{table}[t]
  \centering
  \small
  \caption{Token usage and invocation counts by execution stage.}
  \label{tab:mbpp-token-stages}
  \begin{tabular}{@{}lrrrr@{}}
    \toprule
    Stage & Calls & Input & Output & Total \\
    \midrule
    Code generation      & 8,540 & 2,113,915 & 482,730 & 2,596,645 \\
    EGRCM clarification  & 1,272 &   145,733 &  56,277 &   202,010 \\
    EGRCM summarization  &   424 &   136,077 &  57,725 &   193,802 \\
    Repair               &   344 &    82,373 &  24,287 &   106,660 \\
    \midrule
    Total                &10,580 & 2,478,098 & 621,019 & 3,099,117 \\
    \bottomrule
  \end{tabular}
\end{table}

\paragraph{Why token consumption is high.}
The dominant factor is the fixed ensemble schedule rather than the length or
intrinsic difficulty of MBPP prompts. RE-CODE executes four branches
(EGRCM+PRQE, PRQE, EGRCM, and the unmodified baseline) and samples each branch
five times. Consequently, every problem triggers 20 code-generation calls,
even when an earlier candidate has already passed the tests. Repeated ensemble
generation alone accounts for 2,400,020 tokens (77.4\% of the total), while
all code-generation calls together account for 83.8\%. Relative to one
reference generation per problem (124,510 tokens in total), the complete
pipeline consumes \(24.89\times\) as many tokens.

A second source of overhead is a mismatch between the EGRCM gating heuristic
and the structure of MBPP instances. The gate examines the natural-language
description preceding the function definition, whereas executable examples
and assertions generally occur later in the instance. Thus, the absence of
examples in the inspected text is frequently interpreted as underspecification.
EGRCM was skipped for only 3 of 427 problems; the remaining 424 problems
incurred clarification and summarization. Moreover, the refined requirement
was retained for only one problem and discarded for 423 because its lexical
overlap with the original requirement was below the threshold of 0.35.
Therefore, 395,812 tokens (12.8\% of the total) were spent on EGRCM
clarification and summarization although nearly all resulting refinements were
excluded from code generation.

Prompt augmentation and repair contribute less to the total but remain
non-negligible. PRQE adds four fixed constraints to short prompts, and the
augmented context is resent for every applicable sample; together with the
small amount of accepted EGRCM context, this produces 72,115 additional
tokens (2.3\%). Finally, when all candidates in a branch fail, the repair call
resubmits the task, tests, failure message, and previous program, accounting
for a further 106,660 tokens (3.4\%). The unusually expensive outlier
MBPP/4 (40,276 tokens) exhibits the same pattern: most of its excess
consumption (38,670 tokens) comes from repeated ensemble generation, while
its EGRCM refinement was ultimately rejected.

\paragraph{Limitations and interpretation.}
These results should not be interpreted as showing that MBPP inherently
requires approximately 7.3k tokens per problem. First, the dominant overhead
is induced by our evaluation configuration---four branches, five samples per
branch, and no early stopping---and is therefore an implementation-level
compute--accuracy trade-off. Adaptive branch selection, candidate-level early
stopping, or caching shared prompt prefixes could substantially reduce this
cost without changing the dataset.

Second, the ``reference code generation'' in
Table~\ref{tab:mbpp-token-decomposition} is a counterfactual accounting
baseline formed from the first branch-D generation for each problem; it is
not a separately executed end-to-end baseline. The reported \(24.89\times\)
factor consequently quantifies amplification relative to one generation, not
necessarily the cost ratio against another complete agent framework.

Third, the analysis covers only the 427 completed problems. If incomplete or
failed runs have systematically different token profiles, the mean may be
subject to completion-selection bias. The absolute token counts also depend on
the GPT-4o-mini tokenizer, prompt templates, stopping criteria, and the
specific EGRCM overlap threshold, while the monetary estimate depends on the
prices used at the time of evaluation. Accordingly, comparisons with other
models, datasets, or implementations should report both invocation counts and
input/output tokens rather than extrapolating directly from the dollar cost.

Finally, the fact that only one of 424 refinements was retained indicates that
the current lexical-overlap filter is poorly aligned with EGRCM's behavior on
MBPP. This finding establishes inefficiency under the present configuration,
but it does not isolate whether the primary cause is the gate, the
clarification procedure, or the lexical acceptance criterion. An ablation
that varies these components, together with accuracy and pass-rate
measurements, is required before drawing conclusions about the utility of
EGRCM itself.
